\section*{Chapter \ref{ch:g7:specialisation} --- Specialisation and the Division of Labour}

\begin{solution}{exo:g7:specialisation:1}
Ethiopia $61.0\,\%$, the United States $1.6\,\%$.
\end{solution}

\begin{solution}{exo:g7:specialisation:2}
Alone: 12 chairs per worker a week. Team: $60 \div 4 = 15$ chairs per worker a week.
\end{solution}

\begin{solution}{exo:g7:specialisation:3}
(a) Five steps mean four changes: $4 \times 3 = 12$ minutes per shirt. (b) $40 \times 12
= 480$ minutes, which is $480 \div 60 = 8$ hours.
\end{solution}

\begin{solution}{exo:g7:specialisation:4}
The United States $129.6$, China $8.7$ people per farm worker.
\end{solution}

\begin{solution}{exo:g7:specialisation:5}
(a) $10 + 15 + 5 + 2 \times 5 = 40$ minutes per stool; $480 \div 40 = 12$ stools.
(b) The slowest step takes 12 minutes: $480 \div 12 = 40$ stools. (c) Before: 12 stools
per worker; after: $40 \div 3 \approx 13.3$ stools per worker.
\end{solution}

\begin{solution}{exo:g7:specialisation:6}
(a) No: the village needs 200 pairs a year, less than the 250 a full-time shoemaker
makes. (b) $4 \times 250 = 1\,000$ pairs a year, so a town of $1\,000$ people.
\end{solution}

\begin{solution}{exo:g7:specialisation:7}
The model predicts more distinct jobs: some tasks that were too small to fill a
full-time job (making only soles, only repairing) now are. It gives only the
direction: it does not say how many new kinds of job appear, and nothing says that a
doubled market doubles them.
\end{solution}

\begin{solution}{exo:g7:specialisation:8}
(a) China: $59.6 - 22.2 = 37.4$ percentage points; Korea: $15.5 - 5.2 = 10.3$
percentage points. (b) China: $59.6 \div 22.2 \approx 2.7$ times; Korea: $15.5 \div 5.2
\approx 3.0$ times. Korea's share fell by the larger number of times, although China's
fell by more points.
\end{solution}

\begin{solution}{exo:g7:specialisation:9}
For example: four workers who each make whole toys make 16 a day each, 64 in all;
after the work is split into four steps, with the same tools and hours, the line makes
only 48 a day, 12 per worker.
\end{solution}

\begin{solution}{exo:g7:specialisation:10}
(a) The slowest station takes 5 minutes: $420 \div 5 = 84$ pieces in 7 hours, or 420
minutes. (b) The second station. (c) The slowest station is now the third, at 4
minutes: $420 \div 4 = 105$ pieces. \omterm{def:g7:specialisation:output-per-worker}{Output per worker}: before $84 \div 4 = 21$, after
$105 \div 5 = 21$: unchanged, because the new bottleneck limits the gain.
\end{solution}

\begin{solution}{exo:g7:specialisation:11}
The valley's producers can sell only inside the valley: the market becomes smaller.
Specialists whose work needed outside buyers lose them, so some jobs disappear and
people take up several tasks again (a little farming, a little building). With less
practice, fewer special tools and more time lost between tasks, \omterm{def:g7:specialisation:output-per-worker}{output per worker} falls.
\end{solution}

\begin{solution}{exo:g7:specialisation:12}
(a) In the lower left, far below the curve. (b) One such country does not overturn
the model, but it must be explained. Perhaps its farm workers were counted badly; or
its people left the farms for poorly paid work in towns, or for no work at all, rather
than because farms became more productive; or it \omterm{def:g6:a-connected-world:import}{imports} its food with the income from
one \omterm{def:g6:a-connected-world:export}{export}, such as oil.
\end{solution}

\begin{solution}{pb:g7:specialisation:1}
\textbf{1.} $70 + 50 + 60 + 45 + 3 \times 5 = 240$ minutes, or 4 hours.
\textbf{2.} $480 \div 240 = 2$ bicycles a day, so $2 \times 5 = 10$ a week.
\textbf{3.} $4 \times 10 = 40$ bicycles a week.
\textbf{4.} The frame, at 48 minutes. $480 \div 48 = 10$ bicycles a day, 50 a week.
\textbf{5.} Before: 10 bicycles per worker per week; after: $50 \div 4 = 12.5$. The \omterm{def:g6:reading-economic-charts:rate-of-change}{rate
of change} is $(12.5 - 10) \div 10 = 0.25$, a rise of $25\,\%$.
\textbf{6.} $10 \times 24 = 240$ minutes a day, half of the day.
\textbf{7.} For example: practice (each worker repeats one step all day) and no time
lost changing step; special tools for each step could also help.
\textbf{8.} Frame $40 \div 2 = 20$ minutes; wheels $24 \div 2 = 12$; gears and brakes $32
\div 2 = 16$; check $20 \div 2 = 10$ minutes.
\textbf{9.} The slowest step is the frame, one bicycle every 20 minutes: $480 \div 20 =
24$ a day, $24 \times 5 = 120$ a week.
\textbf{10.} $120 \div 8 = 15$ bicycles per worker per week.
\textbf{11.} It hired $8 - 4 = 4$ workers. \omterm{def:g7:specialisation:output-per-worker}{Output per worker} rose from $12.5$ to 15:
$(15 - 12.5) \div 12.5 = 0.2$, a rise of $20\,\%$.
\textbf{12.} No: it can sell $2 \times 50 = 100$ bicycles a week but builds 120. To sell
them all it needs still more buyers; otherwise it builds only 100, and \omterm{def:g7:specialisation:output-per-worker}{output per worker}
falls back to $100 \div 8 = 12.5$. The size of the market limits how far the \omterm{def:g7:specialisation:division-of-labour}{division
of labour} pays.
\end{solution}
